\scanpage{36}
% archive foliation: 35
i.e.
\diagram{\begin{tikzcd}[column sep=small]
D'^\circ\arrow[d,hook]&\mathbin{\times}&
\underline{\operatorname{Hom}}_{\unsure{k|\mathbf Z}\text{-}add}(A,C')\arrow[r]&C'\\
\underline{\operatorname{Hom}}_{\unsure{k|\mathbf Z}\text{-}add}(A,\mathrm{Ab})^\circ
\end{tikzcd}}
noté
\[
\underline{\operatorname{Hom}}_C(F,G)
\]
----
\par\bigskip\hrule\bigskip
On suppose maintenant qu'on a
une catégorie $\unsure{\mathcal{S}|S}$, une sous-catégorie
\inserted{pl. $C$} de $\unsure{\mathcal{S}|S}_{ab}$ stable par sommes directes
(donc $C$ est additive) et une
sous-catégorie pleine $A$ de $C$ telle
que $\forall$ objet de $C$ est isom. à une
somme directe \deleted{de}\inserted{dans $C$, donc dans $\unsure{\mathcal{S}|S}_{ab}$} d'objets de $A$.
On a alors \deleted{foncteur}
\[
\alpha:\unsure{\mathcal{S}|S}^\circ\longrightarrow\deleted{\unsure{\widehat{\unsure{\mathcal{S}|S}}_{ab}}}
\underline{\operatorname{Hom}}_{add}(C,\mathrm{Ab})
\]
\deleted{\unsure{comme}}
en associant à $\forall X\in\operatorname{Ob}\unsure{\mathcal{S}|S}$
le foncteur $M\longmapsto\operatorname{Hom}_{\unsure{\mathcal{S}|S}}(X,M)$.
Or ce foncteur est connu à équiv.
près par \unsure{la connaissance de} son image
dans $\underline{\operatorname{Hom}}_{\unsure{k|\mathbf Z}\text{-}add}(A,(\mathrm{Ab}))$, donc l'étude
de $\alpha$ est équivalente à celle de
